\section{What the camera can do that light cannot}\label{sec:affordances}

Table~\ref{tab:affordances} collects what follows from the preceding sections
for someone holding the camera. Each line corresponds to a demonstration on
the accompanying site, in which one half of the picture is taken in light and
the other in gravity.

\begin{table}
\begin{tabular}{p{0.26\linewidth}p{0.4\linewidth}p{0.22\linewidth}}
The camera can & Because & Regime \\ \hline
See through anything & nothing absorbs or scatters gravity & near field, thermal \\
Find what is hidden & mass has one sign, so no shield or cloak exists & near field \\
Weigh by looking & the signal is set by mass, not by surface & near field \\
Ignore images & a mirror's reflection or a screen's picture has no mass & all \\
Read inside temperature & hot bodies are transparent to their own glow & thermal \\
Tell materials by their friction & emission follows internal friction & thermal \\
Read spin and its sense & polarization follows the motion of the source & waves \\
See flow inside a sealed body & the magnetic part of the field is made by moving mass & near field \\
See what emits no light & all energy gravitates & near field \\
\end{tabular}
\caption{What a gravitational camera affords, and why.}\label{tab:affordances}
\end{table}

The last line of the table includes light itself. A beam of power $P$ carries
energy and an equal pressure along its length, and to a static observer beside
it both gravitate: it pulls like a line of matter of $2P/c^{3}$ per metre,
twice what its energy alone would suggest \cite{TolmanEhrenfestPodolsky1931}.
An infrared beam that no camera can see from the side is, for this one, a
faint line; for a $100\,\mathrm{kW}$ beam, a change of one part in $10^{22}$
of the pull of the room around it, visible only by subtracting the reading
with the beam off.

The same reasoning tells what the camera cannot do. It cannot read a page,
recognise a face by its colouring, or see a picture on a screen: everything
that light carries by reflection from surfaces is absent. The two kinds of
camera are complementary, and a scene is best understood with both.
